Let \(C/\mathbf C\) be a smooth projective curve, let \(\xi\) be the
good-reduction point of \((C_K)^{\mathrm{an}}\), and let \(B_R\) be the
residual open disc indexed by \(R\in C(\mathbf C)\).  Choose a constant
rational function \(z_R\) which is a local parameter at \(R\), and put
\(r_R=-\log|z_R|\) on the initial radial segment of \(B_R\).  Thus
\[
  \log|g(\zeta_{R,r})|=-\operatorname{ord}_R(g)r
\]
for every \(g\in\mathbf C(C)^\times\) and sufficiently small \(r>0\).

We first prove an obstruction valid for every such constant curve.  If
\(u\) is ordinary smooth near \(\xi\), its local moment presentation has
the form
\[
  u=F(-\log|f_1|,\ldots,-\log|f_m|),
\]
where \(F\) is smooth and the \(f_i\) are invertible analytic functions.
After scaling \(f_i\) by constants of their norms at \(\xi\), their nonzero
reductions are rational functions \(g_i\in\mathbf C(C)^\times\), and on a
smaller neighbourhood the absolute values of the scaled functions and
their reductions agree.  The chain rule therefore shows that the outgoing
slope divisor
\[
 A_u=\sum_R
  \left.\frac{d}{dr}u(\zeta_{R,r})\right|_{r=0+}[R]
   =\sum_i (\partial_iF)(\xi)\operatorname{div}(g_i)
\]
is a finite real linear combination of principal divisors.  In particular
\(\deg A_u=0\).

Let \(D=\sum_Rd_R[R]\) be an integral divisor of degree zero and suppose
\(\mathcal O_{C_K}(D)^{\mathrm{an}}\) has an ordinary smooth psh metric.
For its canonical rational section \(s\), put
\(\lambda=-\log\|s\|\) near \(\xi\), and write
\(A_\lambda=\sum_Ra_R[R]\).  The preceding paragraph makes
\(A_\lambda\) real-principal of degree zero.  On \(B_R\), the section
\(e_R=s/z_R^{d_R}\) is a frame with psh weight
\[
 w_R=-\log\|e_R\|=\lambda-d_Rr_R,
 \qquad
 w_R(\zeta_{R,r})=w_R(\xi)+(a_R-d_R)r+o(r).
\]
Choose a smooth nondecreasing function \(\eta\colon\mathbf R\to[0,1]\)
which is zero on \((-\infty,1]\) and one on \([2,\infty)\).  The function
\(\chi_\varepsilon=\eta(r_R/\varepsilon)\), extended by one at the centre
of \(B_R\), is a smooth nonnegative compactly supported test function.
Positivity of \(d'd''[w_R]\), followed by one-dimensional tropical
integration by parts, gives
\[
  0\leq \langle d'd''[w_R],\chi_\varepsilon\rangle
    =\int_\varepsilon^{2\varepsilon}
       w_R(\zeta_{R,r})\frac{d^2}{dr^2}
       \eta(r/\varepsilon)\,dr
    \longrightarrow -(a_R-d_R).
\]
Hence \(a_R\leq d_R\) for every \(R\).  Since the finite divisors
\(A_\lambda\) and \(D\) both have degree zero, they are equal
coefficientwise.  Thus \(D\) is a finite real linear combination of
integral principal divisors.  This is a finite linear system with rational
coefficients and integral right-hand side.  Consistency over \(\mathbf R\)
implies consistency over \(\mathbf Q\); clearing denominators makes a
positive multiple of \(D\) principal.  Therefore
\([D]\in\operatorname{Pic}^0(C)\) is torsion.

Conversely, if \(mD=\operatorname{div}(f)\), define
\(-\log\|s\|=-(1/m)\log|f|\).  Near \(R\), write
\(f=z_R^{md_R}u_R\), with \(u_R\) a constant-field unit.  Then the weight
of \(s/z_R^{d_R}\) is zero.  These weights give an ordinary smooth
pluriharmonic metric.  We have proved that a constant degree-zero line
bundle has an ordinary smooth psh metric exactly when its divisor class is
torsion.

Return to \(E\).  The point \(P\) has infinite order: for each \(m\geq1\),
the finite morphism \([m]\colon E\to E\) is defined over \(\mathbf Q\), so
all torsion-point coordinates are algebraic over \(\mathbf Q\), whereas
the \(x\)-coordinate of \(P\) is transcendental.  The standard
Riemann--Roch and chord-and-tangent identification
\(E(\mathbf C)\simeq\operatorname{Pic}^0(E)\) sends \(P\) to \([P-O]\).
The obstruction therefore excludes every globally ordinary smooth psh
metric on \(L\).

It remains to prove the local assertion.  In the canonical rational frame
\(s\), define
\[
 \lambda_0=-\log\|s\|_{h_0}=
 \begin{cases}
  r_P&\text{on }B_P\setminus\{P\},\\
  -r_O&\text{on }B_O\setminus\{O\},\\
  0&\text{on all other residual discs and at }\xi.
 \end{cases}
\]
The frame weights of \(s/z_P\) and \(sz_O\) are zero on \(B_P\) and
\(B_O\), respectively, so the metric is continuous at the two centres; it
is continuous at \(\xi\) because only two residual directions are nonzero.
This is the integral model metric from the constant smooth model.  Its
degree on the unique special-fibre curve is zero, hence it is nef.

Away from \(\xi\), its frame weights are constant or logarithms of
invertible analytic functions and are ordinary smooth pluriharmonic.  Fix
\(a>0\), remove the two closed subdiscs \(\{r_P\geq a\}\) and
\(\{r_O\geq a\}\), and call the resulting fixed neighbourhood of \(\xi\)
\(V\).  Then \(s\) is a frame on \(V\) and \(|\lambda_0|\leq a\).

For \(n\geq1\), put \(R_n=(n+1)P\).  These points are distinct from \(P\)
and \(O\), and the elliptic group law supplies \(f_n\in\mathbf C(E)^\times\)
with
\[
 \operatorname{div}(f_n)=(n+1)[P]-n[O]-[R_n].
\]
Choose a smooth convex \(\rho\colon\mathbf R\to\mathbf R\) equal to \(u\)
for \(u\geq-1\) and to \(-2\) for \(u\leq-3\).  Set
\(\mu_n=-(1/n)\log|f_n|\) on \(V\setminus B_{R_n}\), and on \(B_{R_n}\)
set
\[
 \mu_n=\frac1n\rho(-r_{R_n}).
\]
The definitions agree near the boundary.  The function \(\mu_n\) is
ordinary smooth, and its only nonzero Hessian lies in the transition
annulus, where
\[
 d'd''\mu_n=\frac{\rho''(-r_{R_n})}{n}
             d'r_{R_n}\wedge d''r_{R_n}\geq0.
\]
Thus \(\mu_n\) is ordinary smooth psh.  The divisor formula gives
\(\mu_n=(1+1/n)r_P\) on \(B_P\cap V\), \(\mu_n=-r_O\) on
\(B_O\cap V\), zero on all other residual discs except \(B_{R_n}\), and
\(|\mu_n|\leq2/n\) there.  Consequently
\[
  \sup_V|\mu_n-\lambda_0|\leq\frac{\max\{a,2\}}n\longrightarrow0.
\]
The uniform-closure characterization of continuous CLD-psh functions now
makes \(\lambda_0\) psh on \(V\).  Hence every local frame weight of
\(h_0\) is ordinary CLD-psh.  Since \(L\) has no globally ordinary smooth
psh metric, \(h_0\) cannot be a uniform limit of such global metrics.
